\documentclass[10pt,a4paper]{amsart}
\usepackage{amsmath,amssymb,amsthm,graphicx}
% For commutative diagrams we use Xy-pic
\usepackage[all]{xy}

% We use 2cell for 2-commutative diagrams.
\xyoption{2cell}
% Theorem environments.
%
\theoremstyle{plain}
\newtheorem{theorem}[subsection]{Theorem}
\newtheorem{proposition}[subsection]{Proposition}
\newtheorem{lemma}[subsection]{Lemma}

\theoremstyle{definition}
\newtheorem{definition}[subsection]{Definition}
\newtheorem{example}[subsection]{Example}
\newtheorem{exercise}[subsection]{Exercise}
\newtheorem{situation}[subsection]{Situation}

\theoremstyle{remark}
\newtheorem{remark}[subsection]{Remark}
\newtheorem{remarks}[subsection]{Remarks}

\numberwithin{equation}{subsection}

% Macros
%
\def\lim{\mathop{\mathrm{lim}}\nolimits}
\def\colim{\mathop{\mathrm{colim}}\nolimits}
\def\Spec{\mathop{\mathrm{Spec}}}
\def\Hom{\mathop{\mathrm{Hom}}\nolimits}
\def\Ext{\mathop{\mathrm{Ext}}\nolimits}
\def\SheafHom{\mathop{\mathcal{H}\!\mathit{om}}\nolimits}
\def\SheafExt{\mathop{\mathcal{E}\!\mathit{xt}}\nolimits}
\def\Sch{\mathit{Sch}}
\def\Mor{\mathop{\mathrm{Mor}}\nolimits}
\def\Ob{\mathop{\mathrm{Ob}}\nolimits}
\def\Sh{\mathop{\mathit{Sh}}\nolimits}
\def\NL{\mathop{N\!L}\nolimits}
\def\CH{\mathop{\mathrm{CH}}\nolimits}
\def\proetale{{pro\text{-}\acute{e}tale}}
\def\etale{{\acute{e}tale}}
\def\QCoh{\mathit{QCoh}}
\def\Ker{\mathop{\mathrm{Ker}}}
\def\Im{\mathop{\mathrm{Im}}}
\def\Coker{\mathop{\mathrm{Coker}}}
\def\Coim{\mathop{\mathrm{Coim}}}

% Boxtimes
%
\DeclareMathSymbol{\boxtimes}{\mathbin}{AMSa}{"02}

%
% Macros for moduli stacks/spaces
%
\def\QCohstack{\mathcal{QC}\!\mathit{oh}}
\def\Cohstack{\mathcal{C}\!\mathit{oh}}
\def\Spacesstack{\mathcal{S}\!\mathit{paces}}
\def\Quotfunctor{\mathrm{Quot}}
\def\Hilbfunctor{\mathrm{Hilb}}
\def\Curvesstack{\mathcal{C}\!\mathit{urves}}
\def\Polarizedstack{\mathcal{P}\!\mathit{olarized}}
\def\Complexesstack{\mathcal{C}\!\mathit{omplexes}}
% \Pic is the operator that assigns to X its picard group, usage \Pic(X)
% \Picardstack_{X/B} denotes the Picard stack of X over B
% \Picardfunctor_{X/B} denotes the Picard functor of X over B
\def\Pic{\mathop{\mathrm{Pic}}\nolimits}
\def\Picardstack{\mathcal{P}\!\mathit{ic}}
\def\Picardfunctor{\mathrm{Pic}}
\def\Deformationcategory{\mathcal{D}\!\mathit{ef}}
\usepackage{cleveref}
\crefname{originallemma}{original lemma}{original lemmas}
\Crefname{originallemma}{Original lemma}{Original lemmas}
\makeatletter
\newlabel{algebra-lemma-spec-spectral}{{090M}{source}{}{sourcepointer-1}{}}
\newlabel{algebra-lemma-spec-spectral@cref}{{[originallemma][0][]090M}{[1][0][]source}{}{}{}}
\makeatother
\crefname{originalreference}{original reference}{original references}
\Crefname{originalreference}{Original reference}{Original references}
\makeatletter
\newlabel{lemma-gabber-h0}{{09ZF}{source}{}{sourcepointer-2}{}}
\newlabel{lemma-gabber-h0@cref}{{[originalreference][0][]09ZF}{[1][0][]source}{}{}{}}
\makeatother
\crefname{originalreference}{original reference}{original references}
\Crefname{originalreference}{Original reference}{Original references}
\makeatletter
\newlabel{lemma-h0-topological}{{09ZG}{source}{}{sourcepointer-3}{}}
\newlabel{lemma-h0-topological@cref}{{[originalreference][0][]09ZG}{[1][0][]source}{}{}{}}
\makeatother
\crefname{originallemma}{original lemma}{original lemmas}
\Crefname{originallemma}{Original lemma}{Original lemmas}
\makeatletter
\newlabel{more-algebra-lemma-integral-over-henselian-pair}{{09XK}{source}{}{sourcepointer-4}{}}
\newlabel{more-algebra-lemma-integral-over-henselian-pair@cref}{{[originallemma][0][]09XK}{[1][0][]source}{}{}{}}
\makeatother
\crefname{originallemma}{original lemma}{original lemmas}
\Crefname{originallemma}{Original lemma}{Original lemmas}
\makeatletter
\newlabel{more-algebra-lemma-irreducible-henselian-pair-connected}{{09Y6}{source}{}{sourcepointer-5}{}}
\newlabel{more-algebra-lemma-irreducible-henselian-pair-connected@cref}{{[originallemma][0][]09Y6}{[1][0][]source}{}{}{}}
\makeatother
\title{Sections of an etale sheaf over the spectrum of a henselian pair agree with sections over the closed fibre}
\author{The Stacks Project Authors}
\date{stacks-project:etale-cohomology.tex:local-snapshot-20261011; GNU FDL 1.2 or later}
\begin{document}
\maketitle
\noindent\textit{Editorial note.} All delivered mathematical text is the author's own, copied from the pinned Stacks Project chapter file etale-cohomology.tex; no statement or proof step was rewritten, reordered or removed.\par\medskip
\noindent\textit{Editorial note.} The counted proof is the author's complete argument: the connectedness input from the theory of henselian pairs, the Zariski statement for X and Z, its verification after every finite base change, and Gabber's comparison step passing to the etale topology.\par\medskip
\noindent\textit{Editorial note.} The label lemma-h0-henselian-pair names the counted statement, and the proof environment immediately following it in the source is delivered verbatim as its proof.\par\medskip
\noindent\textit{Editorial note.} The counted statement is carried by the author's own lemma environment in the chapter source.\par\medskip
\noindent\textit{Editorial note.} Pre-open gate: python operations/agent-authoring/tools/precard\_gate.py --source etale-cohomology.tex --label lemma-h0-henselian-pair --start 15179 --end 15199 --local-id 3564 returned RESULT=OK (L3564 free, occupancy FREE, public interval clear). Candidate taken from lane-s4/halfB-priority.json.\par\medskip
\noindent\textit{Editorial note.} The statement and proof are copied from the pinned chapter file; the cited inputs are kept as external pointers with their canonical Stacks tags.\par\medskip
\noindent\textit{Allowed standard prerequisites (editorial).}\begin{itemize}
\item The spectrum of a ring is a spectral space, and the trace of a henselian pair on an integral extension is again henselian. Algebra, the lemma that the spectrum of a ring is spectral; More on Algebra, the lemma that an integral extension of a henselian pair is henselian, and the lemma that closed subsets meeting the fibre of a henselian pair stay connected. Supply the connectedness of the trace of the closure of a point with the closed fibre, which is the hypothesis of the topological statement, and the persistence of the henselian property along finite morphisms.
\item The topological statement for the Zariski site. Same chapter, the lemma proving the equality of sections over X and over Z for sheaves on the Zariski site under the connectedness hypothesis. Supplies the equality of sections over the small Zariski site of X and of Z.
\item Gabber's comparison statement for the h-topology. Same chapter, the lemma allowing the passage from the Zariski statement on all finite morphisms to the etale statement. Supplies the final step transferring the verified Zariski statement for all finite morphisms X' -> X to the etale topology.
\end{itemize}\medskip
\ifdefined\hypertarget\hypertarget{sourcepointer-1}{}\fi
\noindent\textit{Source reference.} Original lemma 090M: sources/stacks-chow/tags.txt, source lines 11415--11415. Prerequisite pointer invoked at proof line(s) 15188; it names a standard fact of a sibling chapter that the counted argument uses, and every step of the counted proof built on it is retained verbatim in the delivered body.\par\medskip
\ifdefined\hypertarget\hypertarget{sourcepointer-2}{}\fi
\noindent\textit{Source reference.} Original reference 09ZF: sources/stacks-etale-cohomology/pinned-original-source.tex, source lines 14975--14975. Prerequisite pointer invoked at proof line(s) 15198; it names a standard fact of the same chapter that the counted argument uses, and every step of the counted proof built on it is retained verbatim in the delivered body.\par\medskip
\ifdefined\hypertarget\hypertarget{sourcepointer-3}{}\fi
\noindent\textit{Source reference.} Original reference 09ZG: sources/stacks-etale-cohomology/pinned-original-source.tex, source lines 15099--15099. Prerequisite pointer invoked at proof line(s) 15192; it names a standard fact of the same chapter that the counted argument uses, and every step of the counted proof built on it is retained verbatim in the delivered body.\par\medskip
\ifdefined\hypertarget\hypertarget{sourcepointer-4}{}\fi
\noindent\textit{Source reference.} Original lemma 09XK: sources/stacks-chow/tags.txt, source lines 12533--12533. Prerequisite pointer invoked at proof line(s) 15196; it names a standard fact of a sibling chapter that the counted argument uses, and every step of the counted proof built on it is retained verbatim in the delivered body.\par\medskip
\ifdefined\hypertarget\hypertarget{sourcepointer-5}{}\fi
\noindent\textit{Source reference.} Original lemma 09Y6: sources/stacks-chow/tags.txt, source lines 12554--12554. Prerequisite pointer invoked at proof line(s) 15190; it names a standard fact of a sibling chapter that the counted argument uses, and every step of the counted proof built on it is retained verbatim in the delivered body.\par\medskip
\begin{lemma}
\label{lemma-h0-henselian-pair}
Let $(A, I)$ be a henselian pair. Set $X = \Spec(A)$ and
$Z = \Spec(A/I)$. For any sheaf $\mathcal{F}$ on $X_\etale$
we have $\Gamma(X, \mathcal{F}) = \Gamma(Z, \mathcal{F}|_Z)$.
\end{lemma}

\begin{proof}
Recall that the spectrum of any ring is a spectral space, see
Algebra, Lemma \ref{algebra-lemma-spec-spectral}. By
More on Algebra, Lemma
\ref{more-algebra-lemma-irreducible-henselian-pair-connected}
we see that $\overline{\{x\}} \cap Z$ is connected for every $x \in X$.
By Lemma \ref{lemma-h0-topological} we see that the statement
is true for sheaves on $X_{Zar}$. For any finite morphism $X' \to X$
we have $X' = \Spec(A')$ and $Z \times_X X' = \Spec(A'/IA')$
with $(A', IA')$ a henselian pair, see More on Algebra, Lemma
\ref{more-algebra-lemma-integral-over-henselian-pair}
and we get the same statement for sheaves on $(X')_{Zar}$.
Thus we can apply Lemma \ref{lemma-gabber-h0} to conclude.
\end{proof}

\end{document}
